\begin{abstract}
Supervisors of security operations centres must judge periodically whether monitored
organizations handle security work properly, yet their usual instruments are
self-assessed maturity models and manual record review. We study whether the records
an organization submits can instead be analysed into findings that cite those
records, ranked across organizations, decided by a human supervisor, and bound to that
decision cryptographically. We evaluate this workflow as implemented in SAT-SA, whose
\V{CODE-workers} deterministic, rule-based analytical workers form the evaluated subset
of \V{CODE-registry} registered components, using frozen, hash-verified experiment
bundles. On controlled data the mechanisms behave as designed: every declared finding
family was emitted, validation matched its declared contract in all \V{R01b-expect}
perturbed conditions, all \V{T01-detected} tested mutations of the decided record were
detected, and all \V{O02-runs} injected interruptions recovered. Simple rules are
competitive: fixed-threshold and median-deviation rules tied the fast-closure
detector, and on generated populations that favour SAT-SA by construction a
closure-speed heuristic nearly matched its ranking; stale and contradictory records
were accepted silently. On a public IT incident log, which is not security-operations
data, the risk score was not associated with service-level misses (Spearman
$\rho=\V{X02b-rho-risk}$), whereas slowest median resolution was
($\rho=\V{X02b-rho-slowest}$); a failure analysis attributes this to construct mismatch
and detector saturation. The evidence supports mechanism-level claims, not
effectiveness in real security operations, for which no human study or expert labels
exist.

\keywords{Security operations \and Supervisory analytics \and Incident management
\and Process conformance \and Human-in-the-loop decision support \and Post-quantum
signatures \and Negative results.}
\end{abstract}
